% 实数（综述）
% license CCBYSA3
% type Wiki

本文根据 CC-BY-SA 协议转载翻译自维基百科\href{https://en.wikipedia.org/wiki/Real_number}{相关文章}

在数学中，实数是可以用来度量连续的一维量（如长度、时间或温度）的数。在这里，“连续”意味着数值对之间可以有任意小的差异。[a] 每一个实数都可以通过一个无限小数展开来几乎唯一地表示。[b]

实数在微积分（以及数学的许多其他分支）中是基础性的，尤其体现在其在极限、连续性和导数等经典定义中的作用。[c]

实数集有时被称为 “the reals”（实数），传统上记作黑体 R，通常使用黑板粗体符号 ⁠
$\mathbf{R}$⁠。[1][2] “实数” 这一形容词在 17 世纪由勒内·笛卡尔使用，用以区别实数与虚数（如 −1 的平方根）。[3]

实数包括有理数，例如整数 −5 和分数 $4/3$。其余的实数称为无理数。一些无理数（以及所有有理数）是具有整系数多项式的根，例如平方根 $\sqrt{2} = 1.414...$；这些数被称为代数数。而也存在一些并非如此的实数，例如 $\pi = 3.1415...$；这些数被称为超越数。[3]
\begin{figure}[ht]
\centering
\includegraphics[width=8cm]{./figures/858085ef3576f36f.png}
\caption{实数可以被看作数轴上的所有点。} \label{fig_ShiSu_1}
\end{figure}
实数可以被看作位于一条称为数轴或实数轴上的所有点，其中对应于整数（…，−2，−1，0，1，2，…）的点是等间距的。

上述关于实数的非正式描述并不足以保证涉及实数的定理证明的严谨性。人们逐渐认识到需要一个更精确的定义，并发展出这样的定义，这是 19 世纪数学中的一项重大进展，也是实分析（研究实函数和实数值数列的数学分支）的基础。一个当前的公理化定义是：实数构成唯一的（同构意义下）Dedekind 完备有序域。[d]

实数的其他常见定义还包括：有理数的 Cauchy 数列的等价类、Dedekind 切割以及无限小数展开。这些定义都满足公理化定义，因此是等价的。
\subsection{刻画性质}
实数的完整刻画来自于它们的基本性质，可以总结为：它们构成一个 Dedekind 完备的有序域。这里，“完整刻画” 的意思是，在任意两个 Dedekind 完备有序域之间存在唯一的同构，因此它们的元素具有完全相同的性质。这意味着人们可以运算和操作实数，而无需知道它们是如何被定义的；事实上，在 19 世纪下半叶出现第一个形式化定义之前，数学家和物理学家已经这样做了几个世纪。关于这些形式化定义及其等价性的证明，详见实数的构造。
\subsection{算术}
实数构成一个有序域。直观上，这意味着初等算术的方法和规则对实数成立。更精确地说，实数上定义了两个二元运算（加法与乘法）以及一个全序关系，并且满足以下性质：
\begin{itemize}
\item 加法：两个实数 $a$ 和 $b$ 的加法产生一个实数，记作$a+b$，它是 $a$ 和 $b$ 的和。
\item 乘法：两个实数 $a$ 和 $b$ 的乘法产生一个实数，记作$ab,\quad a\cdot b,\quad a\times b$，它是 $a$ 和 $b$ 的积。
\item 交换律：加法和乘法都是交换的，即$a+b=b+a,\quad ab=ba$对任意实数 $a, b$ 都成立。
\item 结合律：加法和乘法都是结合的，即$(a+b)+c=a+(b+c),\quad (ab)c=a(bc)$对任意实数 $a, b, c$ 都成立，因此在这两种运算中括号可以省略。
\item 分配律：乘法对加法分配，即$a(b+c)=ab+ac$对任意实数 $a, b, c$ 都成立。
\item 加法单位元：存在一个称为零的实数，记作 $0$，它是加法的单位元，即$a+0=a$
对任意实数 $a$ 都成立。
\item 乘法单位元：存在一个实数，记作 $1$，它是乘法的单位元，即$a\times 1=a$
对任意实数 $a$ 都成立。
\item 加法逆元：每个实数 $a$ 都有一个加法逆元，记作 $-a$，即$a+(-a)=0$对任意实数 $a$ 都成立。
\item 乘法逆元：每个非零实数 $a$ 都有一个乘法逆元，记作$a^{-1}\quad \text{或}\quad \tfrac{1}{a}$，即$aa^{-1}=1$对任意非零实数 $a$ 都成立。
\item 全序：实数上的全序关系记作 $a<b$。它满足两个性质：对于任意两个实数 $a, b$，三者中恰有一个成立：$a<b,\quad a=b,\quad b<a$。若$a<b$且$b<c$，则有 $a<c$。
\item 序与运算的相容性：若 $a<b$，则$a+c<b+c$对任意实数 $c$ 都成立；若 $0<a$且$0<b$，则$0<ab$。
\end{itemize}
从上述性质中可以推导出许多其他性质。特别是：
\begin{itemize}
\item $0 \cdot a = 0$对任意实数 $a$ 都成立。
\item $0 < 1$
\item $0 < a^{2}$对任意非零实数 $a$ 都成立。
\end{itemize}
\subsubsection{辅助运算}
还有一些常用的运算，可以由前述性质推导而来：
\begin{itemize}
\item 减法：两个实数$a$和$b$的减法等于$a$与$b$的加法逆元$-b$的和，即
  $$
  a-b=a+(-b)~
  $$
\item 除法：实数$a$除以一个非零实数$b$，记作$\tfrac{a}{b}, \quad a/b$
，定义为$a$与$b$的乘法逆元的积，即
  $$
  \tfrac{a}{b}=ab^{-1}~
  $$
\item 绝对值：实数$a$的绝对值，记作$|a|$，表示它与零的距离，定义为
  $$
  |a|=\max(a,-a)~
  $$
\end{itemize}
